% Section 9: branch-and-bound for MILP.
\section{Branch-and-bound for mixed-integer programmes}\label{sec:mip}

Every node of the search is the LP relaxation \eqref{eq:simplex-form} with some integer bounds
tightened. A branching change leaves the parent's optimal basis dual feasible, so the dual simplex
re-optimises a child in a few pivots. This is why MILP codes are built on the dual simplex, and
\file{mip/bnb.py} is built on the warm-started \code{SimplexLP} of \cref{sec:simplex}. Its components
follow Achterberg's thesis~\cite{achterberg2007} and the branching study of Achterberg, Koch and
Martin~\cite{achterberg2005}.

\subsection{The search}
\paragraph{Node selection: best bound with plunging}
Open nodes wait in a binary heap keyed by
\begin{equation}\label{eq:heapkey}
  \bigl(\ \underline{z}_{\text{parent}},\ -\text{depth},\ \text{counter}\ \bigr),
\end{equation}
so the node with the smallest bound comes first, and among equal bounds the \emph{deepest}
(\cref{sec:heap}). After a node is branched on, the child that the fractional part favours ($x_j$
rounded up if $f_j\ge0.5$, else down) is solved at once, and the other goes to the heap. The search
keeps diving into the solved child while its bound stays within $10\%$ of the distance between the
best open bound and the cutoff:
\begin{equation}\label{eq:plunge}
  \underline z_{\text{child}} \le \underline z_{\text{open}} + 0.1\,\bigl|z_{\text{cutoff}}-\underline z_{\text{open}}\bigr| .
\end{equation}
Before an incumbent exists the right-hand side is infinite, so the search dives depth-first, which
finds incumbents early. Once one exists it returns to best-first search, which proves the bound.
A child that fails \eqref{eq:plunge} is queued with its own LP bound and basis.

\paragraph{Pruning and bound rounding}
A node is pruned when its LP is infeasible or its rounded bound reaches the cutoff
\begin{equation}\label{eq:cutoff}
  z_{\text{cutoff}} = \begin{cases} z_{\text{inc}} - 1 + 10^{-6} & \text{if the objective is integral on integer points},\\
  z_{\text{inc}} - \max(10^{-9}, 10^{-9}|z_{\text{inc}}|) & \text{otherwise.}\end{cases}
\end{equation}
The objective is integral when every continuous column has zero cost and all costs and $c_0$ are
integers. Node bounds are then rounded up, $\lceil \underline z - 10^{-6}\rceil$, which often closes
the last unit of gap without enumeration. The search stops when the relative gap
$(z_{\text{inc}}-\underline z)/\max(1,|z_{\text{inc}}|)$ is at most $10^{-4}$, the default of HiGHS
and SCIP.

\begin{figure}[tbp]
\centering
\begin{forest}
  for tree={
    draw=navy, rounded corners=2pt, align=center, font=\sffamily\scriptsize,
    inner sep=2pt, minimum width=11mm, l sep=7mm, s sep=9mm, edge={->,draw=ink,line width=0.5pt},
    fill=white,
  }
  [{root\\$\underline z{=}100.3$},fill=navy!12
    [{$x_3\le2$\\$100.6$},edge label={node[midway,left,lbl]{\#1 plunge}},fill=saffron!18,draw=saffron
      [{$x_7\le0$\\$100.8$},edge label={node[midway,left,lbl]{\#2}},fill=saffron!18,draw=saffron
        [{$x_1\le4$\\infeasible},fill=card,draw=muted,edge label={node[midway,left,lbl]{\#3}}]
        [{$x_1\ge5$\\$102$ integral},fill=igreen!12,draw=igreen,edge label={node[midway,right,lbl]{\#4}}]
      ]
      [{$x_7\ge1$\\$100.9$},fill=sky!12,draw=sky,edge label={node[midway,right,lbl]{\#5 from heap}}
        [{$x_2\le1$\\\textbf{101} integral},fill=igreen!25,draw=igreen,edge label={node[midway,left,lbl]{\#6}}]
        [{$x_2\ge2$\\not solved},fill=card,draw=muted,edge label={node[midway,right,lbl]{pruned}}]
      ]
    ]
    [{$x_3\ge3$\\not solved},fill=card,draw=muted,edge label={node[midway,right,lbl]{pruned}}]
  ]
\end{forest}
\caption[Branch-and-bound tree]{An illustrative branch-and-bound tree (made-up numbers, integral
objective). Orange: the plunge, which dives into the favoured child after each branching and finds
the incumbent 102 at node~\#4. The cutoff is then $102-1+10^{-6}$. The open nodes $x_7\ge1$ and
$x_3\ge3$ both have the rounded key $\lceil100.6\rceil=\lceil100.3\rceil=101$; the depth tie-break
of \eqref{eq:heapkey} pops the deeper one (\#5), whose child \#6 is the optimum 101. The new cutoff
$100+10^{-6}$ prunes the two remaining nodes without solving their LPs, and the root bound
$\lceil100.3\rceil=101$ closes the gap.}
\label{fig:bnb}
\end{figure}

\paragraph{Branching: pseudocosts with reliability}
The candidate score is the product rule
\begin{equation}\label{eq:score}
  \text{score}_j = \max\bigl(\psi^-_j f_j,\ 10^{-6}\bigr)\cdot\max\bigl(\psi^+_j(1-f_j),\ 10^{-6}\bigr),
\end{equation}
where $f_j$ is the fractional part of $x_j$ and $\psi^\pm_j$ are the pseudocosts: the average gain in
the LP objective per unit of change, recorded from every solved child. Unobserved directions use the
average over all columns. A candidate is \emph{unreliable} while it has fewer than 4 observations in
either direction. Up to 8 of the best-scored unreliable candidates are then strong-branched, both
children solved by the dual simplex for at most 60 iterations from the current basis, which also
records pseudocosts~\cite{achterberg2005}. If both children are infeasible the node is pruned. If one
is, the candidate gets a top score, because branching on it fixes a variable. An integral child LP
found during strong branching is offered as an incumbent.

\subsection{Domain propagation}\label{sec:propagation}
At every node the bounds are propagated over the rows~\cite{savelsbergh1994,achterberg2007}
(\file{mip/propagate.py}). For a row $r^l\le\sum_j a_jx_j\le r^u$, the minimum and maximum activities
of the row without $x_j$ imply
\begin{equation}\label{eq:propagate}
  a_j>0:\quad x_j\le \frac{r^u-\underline{\text{act}}_{-j}}{a_j},\qquad x_j\ge\frac{r^l-\overline{\text{act}}_{-j}}{a_j},
\end{equation}
with the inequalities reversed for $a_j<0$. Infinite contributions are counted, so a row with exactly
one infinite contributor can still bound that one variable. Integer bounds are rounded inward
($\lfloor u+10^{-6}\rfloor$, $\lceil \ell-10^{-6}\rceil$). A continuous bound is accepted only if it
improves by more than $10^{-3}\max(1,|\text{bound}|)$, so round-off cannot creep into the bounds.
An empty domain, or an activity range that misses the row, proves the node infeasible. Up to 5 passes
run per node. Bounds found for continuous columns guide the search but are not sent to the LP.

\subsection{Primal heuristics}
\begin{itemize}
  \item \textbf{Simple rounding} at every node: round all integer columns of the LP solution and
        accept the point if it is feasible.
  \item \textbf{Fix-and-solve} at the root and every 100 nodes: fix every integer to its rounded
        LP value, propagate, and re-solve the LP for the continuous columns (at most 5\,000
        dual iterations).
  \item \textbf{Fractional diving}~\cite{berthold2006} at the root when no incumbent exists, and every
        50 nodes while none does: repeatedly fix the least fractional integer to its nearest value,
        propagate and re-solve, until the LP is integral, infeasible or worse than the cutoff. The
        LP budget is $\max(1000,\ 2\times$ the LP iterations spent so far$)$.
\end{itemize}

\subsection{Gomory mixed-integer cuts}\label{sec:gmi}
At the root, up to 20 rounds of Gomory mixed-integer (GMI) cuts~\cite{gomory1960,balas1996} tighten
the relaxation (\file{mip/cuts.py}). Take an optimal tableau row whose basic variable is an integer
column with fractional value. Write each nonbasic variable as its distance from the bound it sits
at, $t_k\ge0$, so that $x_B + \sum_k a_k t_k = b$ with $f_0=\fr(b)\in(0,1)$. Every integer-feasible
point then satisfies
\begin{equation}\label{eq:gmi}
  \sum_{\substack{k\ \text{int}\\ f_k\le f_0}}\frac{f_k}{f_0}t_k
  +\sum_{\substack{k\ \text{int}\\ f_k> f_0}}\frac{1-f_k}{1-f_0}t_k
  +\sum_{\substack{k\ \text{cont}\\ a_k>0}}\frac{a_k}{f_0}t_k
  +\sum_{\substack{k\ \text{cont}\\ a_k<0}}\frac{-a_k}{1-f_0}t_k \;\ge\; 1,
  \qquad f_k=\fr(a_k).
\end{equation}
The LP point has $t=0$ and violates the cut by exactly~1. A nonbasic variable counts as integer when its
column is integer and its bound is integral. Logical variables are expanded into the structural
columns of their row, and the cut is mapped back to the unscaled $x$ space.

\paragraph{Numerical safeguards}
Following the practice of Cook, Dash, Fukasawa and Goycoolea~\cite{cook2009}, a cut is generated only
from a row with $0.01<f_0<0.99$. It is skipped when a free nonbasic variable appears in the row or a
needed bound is infinite. Coefficients below $10^{-9}$ of the largest are removed by relaxing against
the variable's bound, which keeps the cut valid. The cut is rejected if its coefficient ratio exceeds
$10^6$ or its efficacy (violation divided by $\|g\|_2$) is below $10^{-5}$, and the survivors are
scaled so that the largest coefficient is~1. A regression test checks that every cut generated on
\code{p0033} keeps the known optimal integer point feasible.

\paragraph{The cut loop}
Each round adds up to $\max(20,\min(200,m/2))$ cuts, most fractional rows first, and re-solves the LP.
The loop stops when the LP is integral, no cut is found, three rounds together move the bound by less
than $0.1\%$, or the cuts have grown the model beyond $4m+1000$ rows. Cuts that are slack at the final
root point are then dropped, because they would only slow down the node LPs.

\begin{algorithm}[t]
\caption{Branch-and-bound with plunging (\file{mip/bnb.py}, simplified)}\label{alg:bnb}
\Input{MILP, gap tolerance $10^{-4}$, time limit}
propagate at the root; solve the root LP; GMI cut rounds; rounding, fix-and-solve, diving\;
$\text{current}\leftarrow$ root (LP solved); heap $\leftarrow\emptyset$\;
\While{time remains}{
  \If{$\text{current}=\emptyset$}{
    pop nodes whose bound $\ge z_{\text{cutoff}}$; \lIf{heap empty}{\textbf{break}}
    pop node with the smallest $(\underline z,-\text{depth})$; propagate; warm-start the dual simplex from its stored basis\;
    record the pseudocost of the branching that created it\;
  }
  \lIf{LP infeasible or $\lceil\underline z\rceil\ge z_{\text{cutoff}}$}{\textbf{continue}}
  \lIf{LP solution integral}{update incumbent; \textbf{continue}}
  simple rounding; periodic diving and fix-and-solve\;
  choose $j$ by reliability branching \eqref{eq:score}\;
  create children $x_j\le\lfloor x_j\rfloor$ and $x_j\ge\lceil x_j\rceil$ with the parent's basis\;
  push the less favoured child; solve the favoured one\;
  \leIf{its bound satisfies \eqref{eq:plunge}}{$\text{current}\leftarrow$ that child}{push it; $\text{current}\leftarrow\emptyset$}
  \lIf{relative gap $\le10^{-4}$}{\textbf{break}}
}
\end{algorithm}
